\documentclass[12pt, a4paper]{article}

% Pakiety
\usepackage[utf8]{inputenc}
\usepackage[T1]{fontenc}
\usepackage[polish]{babel}
\usepackage{amsmath, amssymb}
\usepackage{graphicx}
\usepackage{float}
\usepackage{geometry}
\usepackage{hyperref}
\usepackage{booktabs}
\usepackage{caption}
\usepackage{listings}
\usepackage{xcolor}
\usepackage{parskip}

\geometry{
  a4paper,
  left=2.5cm,
  right=2.5cm,
  top=2.5cm,
  bottom=2.5cm
}

% Zakaz dzielenia wyrazów
\hyphenpenalty=10000
\exhyphenpenalty=10000
\emergencystretch=3em

\hypersetup{
  colorlinks=true,
  linkcolor=blue,
  urlcolor=blue,
  citecolor=blue
}

% Zmiana nazwy float figure: Rysunek -> Wykres
\addto\captionspolish{\renewcommand{\figurename}{Wykres}}

% Styl listingów
\lstset{
  language=Python,
  basicstyle=\ttfamily\small,
  keywordstyle=\color{blue}\bfseries,
  commentstyle=\color{gray}\itshape,
  stringstyle=\color{orange!80!black},
  numbers=left,
  numberstyle=\tiny\color{gray},
  numbersep=8pt,
  frame=single,
  breaklines=true,
  breakatwhitespace=true,
  captionpos=b,
  showstringspaces=false,
  columns=flexible,
  keepspaces=true,
  inputencoding=utf8,
  extendedchars=true,
  literate={ą}{{\k{a}}}1 {ę}{{\k{e}}}1 {ó}{{\'{o}}}1
            {ś}{{\'{s}}}1 {ł}{{\l{}}}1 {ź}{{\'{z}}}1
            {ż}{{\.z}}1 {ć}{{\'{c}}}1 {ń}{{\'{n}}}1
            {Ą}{{\k{A}}}1 {Ę}{{\k{E}}}1 {Ó}{{\'{O}}}1
            {Ś}{{\'{S}}}1 {Ł}{{L}}1,
}

% Strona tytułowa
\title{
  \textbf{Metody Obliczeniowe w Nauce i Technice} \\[0.5em]
  \large Laboratorium 9 \textendash{} Równania różniczkowe zwyczajne (część I)
}
\author{Mateusz Stoch}
\date{13 maja 2026}

\begin{document}

\maketitle

% ============================================================
\section{Zadanie 1 \textendash{} Sprowadzanie do układu równań pierwszego rzędu}
% ============================================================
\textbf{Polecenie:} Przedstaw każde z poniższych równań różniczkowych zwyczajnych jako równoważny układ równań pierwszego rzędu.

Rozwiązanie równań wyższego rzędu z wykorzystaniem standardowych metod numerycznych wymaga uprzedniego przekształcenia ich do układu równań pierwszego rzędu poprzez wprowadzenie zmiennych pomocniczych reprezentujących kolejne pochodne.

\textbf{(a) Równanie Van der Pol'a:} $y'' = y'(1 - y^2) - y$ \\
Wprowadzono zmienne $u_1 = y$ oraz $u_2 = y'$.
\[
  \begin{cases}
    u_1' = u_2 \\
    u_2' = u_2(1 - u_1^2) - u_1
  \end{cases}
\]

\textbf{(b) Równanie Blasiusa:} $y''' = -yy''$ \\
Wprowadzono zmienne $u_1 = y$, $u_2 = y'$ oraz $u_3 = y''$.
\[
  \begin{cases}
    u_1' = u_2 \\
    u_2' = u_3 \\
    u_3' = -u_1 u_3
  \end{cases}
\]

\textbf{(c) Problem dwóch ciał:} 
\[ y''_1 = -\frac{GMy_1}{(y_1^2 + y_2^2)^{3/2}}, \quad y''_2 = -\frac{GMy_2}{(y_1^2 + y_2^2)^{3/2}} \]
Wprowadzono zmienne $u_1 = y_1$, $u_2 = y_1'$, $u_3 = y_2$ oraz $u_4 = y_2'$. Oznaczając $r = (u_1^2 + u_3^2)^{3/2}$, uzyskano:
\[
  \begin{cases}
    u_1' = u_2 \\
    u_2' = -GM u_1 / r \\
    u_3' = u_4 \\
    u_4' = -GM u_3 / r
  \end{cases}
\]

% ============================================================
\section{Zadanie 2 \textendash{} Autonomizacja problemu}
% ============================================================
\textbf{Polecenie:} Przekształć poniższy problem początkowy do autonomicznego problemu początkowego:
\[
  \begin{cases}
    y_1' = y_1/t + y_2 t \\
    y_2' = t(y_2^2 - 1)/y_1
  \end{cases} \quad y_1(1) = 1, \quad y_2(1) = 0
\]

Aby wyeliminować jawną zależność od zmiennej czasu $t$, wprowadza się dodatkową zmienną stanu $u_3 = t$, dla której $u_3' = 1$. Zastępując $y_1$ i $y_2$ przez $u_1$ i $u_2$, układ autonomiczny przyjmuje postać:
\[
  \begin{cases}
    u_1' = u_1/u_3 + u_2 u_3 \\
    u_2' = u_3(u_2^2 - 1)/u_1 \\
    u_3' = 1
  \end{cases} \quad u_1(1)=1, \quad u_2(1)=0, \quad u_3(1)=1
\]

% ============================================================
\section{Zadanie 3 \textendash{} Weryfikacja rozwiązania}
% ============================================================
\textbf{Polecenie:} Dany jest problem początkowy: $y' = \sqrt{1 - y}$, $y(0) = 0$. Pokaż, że funkcja $y(t) = t(4 - t)/4$ spełnia równanie i warunek początkowy oraz wyznacz dziedzinę.

Warunek początkowy: $y(0) = \frac{0 \cdot 4}{4} = 0$. \\
Pochodna badanej funkcji wynosi: $y'(t) = 1 - \frac{t}{2}$. \\
Podstawienie do równania różniczkowego daje:
\[
  \sqrt{1 - y(t)} = \sqrt{\frac{4 - 4t + t^2}{4}} = \sqrt{\left(\frac{2-t}{2}\right)^2} = \left| 1 - \frac{t}{2} \right|
\]
Równość $y' = \sqrt{1-y}$ zachodzi pod warunkiem, że wynik pierwiastkowania jest nieujemny ($1 - t/2 \geq 0$). Wobec warunku $t \geq 0$, wyznaczona dziedzina to $t \in [0, 2]$.

% ============================================================
\section{Zadanie 4 \textendash{} Analiza i numeryka metody Eulera}
% ============================================================
\textbf{Polecenie:} Dane jest równanie $y' = -5y$ z warunkiem początkowym $y(0) = 1$. Równanie rozwiązuje się numerycznie z krokiem $h = 0.5$. Rozpatrz poniższe podpunkty.

\begin{itemize}
  \item[\textbf{(a)}] \textbf{Analityczna stabilność:} Po wprowadzeniu zaburzenia $\varepsilon$ do warunku początkowego, rozwiązanie przyjmuje postać $y(t) = (1+\varepsilon)e^{-5t}$. Różnica między rozwiązaniem zaburzonym a dokładnym wynosi $\delta y(t) = \varepsilon e^{-5t}$. Ponieważ $\lambda = -5 < 0$, błąd ten zanika w czasie ($\lim_{t \to \infty} \varepsilon e^{-5t} = 0$), co oznacza, że problem jest asymptotycznie stabilny.

  \item[\textbf{(b)}] \textbf{Zbieżność metody Eulera:} Wzór na jawną metodę Eulera to $y_{n+1} = y_n + h(-5y_n) = y_n(1-5h)$. Po $n$ iteracjach startując z $y_0 = 1$ otrzymuje się $y_n = (1-5h)^n$. Przekształcając to dla ustalonego punktu $t = nh$, uzyskuje się:
  \[
    y_n = (1-5h)^{t/h} = \left[ (1-5h)^{\frac{1}{5h}} \right]^{-5t}
  \]
  Przy $h \to 0$, wewnętrzne wyrażenie z definicji dąży do $e^{-1}$, a w efekcie cały wynik dąży do wartości dokładnego rozwiązania $e^{-5t}$, co dowodzi zbieżności metody.

  \item[\textbf{(c)}] \textbf{Numeryczna stabilność metody jawnej:} Metoda jest stabilna, gdy tzw. czynnik wzmocnienia $|1 + h\lambda| \leq 1$. Obliczenia dla zadanego kroku dają:
  \[
    |1 + 0.5 \cdot (-5)| = |1 - 2.5| = |-1.5| = 1.5 > 1
  \]
  Przekroczenie jedności oznacza, że przy kroku $h=0.5$ metoda staje się całkowicie niestabilna i obarczona rosnącym błędem.

  \item[\textbf{(d)}] \textbf{Obliczenia numeryczne dla metody jawnej:} Zgodnie z wyprowadzonym wyżej wzorem, w węźle $t=0.5$ po jednym kroku metoda daje przybliżenie:
  \[
    y_1 = y_0(1-5 \cdot 0.5) = 1 \cdot (-1.5) = -1.5
  \]
  Ze względu na niestabilność, wynik ten jest ujemny i kompletnie błędny względem precyzyjnej wartości $y(0.5) = e^{-2.5} \approx 0.082$.

  \item[\textbf{(e)}] \textbf{Stabilność niejawnej metody Eulera:} W metodzie niejawnej kolejny węzeł wyznacza się z użyciem wartości funkcji w nowym punkcie: $y_{n+1} = y_n - 5h y_{n+1}$. Po przekształceniach daje to $y_{n+1} = \frac{y_n}{1 + 5h}$. Współczynnik wzmocnienia wyliczony dla tego przypadku to:
  \[
    \left|\frac{1}{1+5h}\right| = \left|\frac{1}{1+2.5}\right| = \frac{1}{3.5} \approx 0.286 < 1
  \]
  Otrzymana wartość $< 1$ gwarantuje metodyczną A-stabilność dla zastosowanego kroku $h$.

  \item[\textbf{(f)}] \textbf{Obliczenia numeryczne dla metody niejawnej:} Podstawiając zadany krok do wyprowadzonego wariantu niejawnego:
  \[
    y_1 = \frac{y_0}{1+5h} = \frac{1}{1+2.5} = \frac{1}{3.5} \approx 0.286
  \]
  Mimo ograniczonej precyzji ze względu na wielki krok $h=0.5$ (odchylenie od $0.082$), wyliczony wynik pozostaje stabilny numerycznie i przyjmuje sensowne wartości dodatnie.

  \item[\textbf{(g)}] \textbf{Maksymalna dopuszczalna wartość kroku $h$:} Globalny błąd algorytmu opiera się na szacowaniu $|y_n - y(t)| \approx \frac{h}{2} \max |y''(\xi)| t$. Skoro $y(t) = e^{-5t}$, to $y''(t) = 25e^{-5t}$. Jej maksymalna wartość w przedziale $[0, 0.5]$ znajduje się w $t=0$ i wynosi $25$. Rozwiązując nierówność dla błędu $0.001$:
  \[
    \frac{h}{2} \cdot 25 \cdot 0.5 \leq 0.001 \implies 6.25h \leq 0.001 \implies h \leq 1.6 \times 10^{-4}
  \]
  Aby dotrzeć numerycznie do $t=0.5$ przy tak małym kroku, należałoby wykonać $n = 0.5 / (1.6 \times 10^{-4}) \approx 3125$ kroków.

  \item[\textbf{(h)}] \textbf{Metoda wyznaczania $y_{n+1}$:} Iteracja prosta dla niejawnego Eulera przyjmuje postać $\varphi(y) = y_n - 5hy$. Rozwiązanie tego algorytmu jest zbieżne z twierdzenia Banacha tylko gdy $|\varphi'(y)| < 1$, czyli w tym wypadku $|-5h| < 1 \implies h < 0.2$. Z tego powodu dla narzuconego problemu kroku $h=0.5$ iteracja prosta od razu by się rozbiegła. W związku z tym, użycie w to miejsce metody Newtona (która nie ma takiego uciążliwego ograniczenia na rozmiar kroku) jest jedynym prawidłowym wyjściem.
\end{itemize}

% ============================================================
\section{Zadanie 5 \textendash{} Stabilność układu równań}
% ============================================================
\textbf{Polecenie:} Dany jest układ równań $y'_1 = -2y_1 + y_2$, $y'_2 = -y_1 - 2y_2$. Wyznacz dopuszczalne wartości kroku $h$ dla stabilności metody Eulera.

\begin{itemize}
  \item \textbf{Zapis macierzowy:} Rozwiązanie stabilności układu równań liniowych sprowadza się do analizy wartości własnych jego macierzy współczynników. Zapisując podany układ w postaci $\mathbf{y}' = A \mathbf{y}$, otrzymuje się:
  \[
    A = \begin{pmatrix} -2 & 1 \\ -1 & -2 \end{pmatrix}
  \]

  \item \textbf{Wartości własne:} Następnie wyznaczane są wartości własne z równania charakterystycznego $\det(A - \lambda I) = 0$:
  \[
    \det \begin{pmatrix} -2-\lambda & 1 \\ -1 & -2-\lambda \end{pmatrix} = (-2-\lambda)^2 - (-1) \cdot 1 = (\lambda+2)^2 + 1 = 0
  \]
  Rozwiązaniem tego równania w dziedzinie liczb zespolonych są sprzężone wartości własne: $\lambda_{1,2} = -2 \pm i$.

  \item \textbf{Kryterium stabilności:} Dla jawnej metody Eulera obszar stabilności bezwzględnej zdefiniowany jest jako zbiór takich wartości zespolonych $h\lambda$, dla których czynnik wzmocnienia leży wewnątrz okręgu jednostkowego. Warunek ten wyraża się nierównością $|1 + h\lambda| \leq 1$. 

  \item \textbf{Rozwiązanie nierówności:} Podstawiając jedną z wyliczonych wartości $\lambda = -2 + i$, otrzymuje się liczbę zespoloną z częścią rzeczywistą $(1-2h)$ oraz urojoną $h$. Kwadrat modułu tej liczby nie może przekroczyć $1$:
  \[
    | (1-2h) + ih |^2 \leq 1 \implies (1-2h)^2 + h^2 \leq 1
  \]
  Po rozwinięciu wzoru skróconego mnożenia powstaje równanie:
  \[
    1 - 4h + 4h^2 + h^2 \leq 1 \implies 5h^2 - 4h \leq 0
  \]

  \item \textbf{Dopuszczalny krok:} Wyciągając zmienną $h$ przed nawias $h(5h - 4) \leq 0$ oraz zakładając, że krok czasowy musi być z natury wartością dodatnią, uzyskuje się ostateczne kryterium. Metoda Eulera dla tego konkretnego układu pozostaje stabilna numerycznie dla rozmiaru kroku ograniczonego przez:
  \[
    h \leq \frac{4}{5} \implies h \in (0, 0.8]
  \]
\end{itemize}

% ============================================================
\section{Zadanie 6 \textendash{} Rząd zbieżności metody Eulera}
% ============================================================
\textbf{Polecenie:} Dany jest problem $y' = \alpha t^{\alpha-1}$, $y(0) = 0$. Wyznacz numerycznie wartości błędu oraz empiryczny rząd zbieżności dla $\alpha \in \{2.5, 1.5, 1.1\}$ i kroków $h \in \{0.2, 0.1, 0.05\}$. Wyjaśnij otrzymane wyniki.

\begin{itemize}
  \item \textbf{Implementacja algorytmu:} Do rozwiązania zadania napisano dedykowaną funkcję wykonującą symulację dla podanych parametrów kroku. Kod skróconej implementacji jawnej metody Eulera prezentuje się następująco:
\begin{lstlisting}[caption={Skrócony kod rozwiązania}]
def euler_solve(alpha, h, t_end=1.0):
    t, y = 0.0, 0.0
    ts, ys = [t], [y]
    while t < t_end - h/2:
        dydt = alpha * t**(alpha - 1.0) if t > 0 else 0.0
        y = y + h * dydt
        t = t + h
        ts.append(t), ys.append(y)
    return np.array(ts), np.array(ys)
\end{lstlisting}

  \item \textbf{Tabela błędów:} Poniższe zestawienie pokazuje, jak zachowywał się błąd bezwzględny dla poszczególnych badanych parametrów $\alpha$ w węźle $t=1$ i dla zmniejszających się rozmiarów kroku:
\begin{center}
\begin{tabular}{r|r|r|r}
\toprule
$\alpha$ & Błąd dla $h=0.2$ & Błąd dla $h=0.1$ & Błąd dla $h=0.05$ \\
\midrule
$2.5$ & $0{,}2386$ & $0{,}1220$ & $0{,}0617$ \\
$1.5$ & $0{,}1753$ & $0{,}0842$ & $0{,}0408$ \\
$1.1$ & $0{,}1877$ & $0{,}0913$ & $0{,}0444$ \\
\bottomrule
\end{tabular}
\end{center}

  \item \textbf{Obliczony rząd zbieżności:} Rząd wyliczono empirycznie z powyższego zestawienia błędów. Wyniki dla poszczególnych stałych potęg wynoszą odpowiednio: dla $\alpha=2.5 \implies p \approx 0.98$, dla $\alpha=1.5 \implies p \approx 1.04$, dla $\alpha=1.1 \implies p \approx 1.03$. 
  
  \item \textbf{Analiza wyników:} W każdym testowanym wariancie wyliczony rząd oscyluje wokół $p=1$. Wynika to z faktu, że metoda Eulera ze swojej natury charakteryzuje się błędem o teoretycznej złożoności $O(h^p)$, gdzie dla równań z potęgą $p = \min(1, \alpha)$. Przy badanych parametrach $\alpha > 1$, prawostronna pochodna wykazuje wymaganą gładkość dla $t=0$, potwierdzając numerycznie i analitycznie liniowy rząd zbieżności zastosowanego algorytmu.

  \item \textbf{Wizualizacja (Wykres 1):} Wymienione zbieżności i wyniki błędów przedstawiono dodatkowo w formie graficznej na poniższym wykresie.
  
  \begin{figure}[H]
      \centering
      \includegraphics[width=0.9\textwidth]{wykres1.png}
      \caption{Wykres nr 1 - Badanie rzędu zbieżności metody Eulera}
      \label{fig:wykres1}
  \end{figure}
\end{itemize}

% ============================================================
\section{Podsumowanie}
% ============================================================
Na podstawie przeprowadzonych w laboratorium zadań wyciągnięto następujące wnioski:
\begin{itemize}
    \item Jawna metoda Eulera cechuje się kiepską stabilnością i wymaga narzucenia bardzo małych kroków $h$, co często jest nieefektywne.
    \item Niejawna metoda Eulera zachowuje świetną stabilność nawet dla dużych kroków, wymaga jednak zastosowania w każdej iteracji dodatkowych obliczeń.
    \item Metoda Eulera posiada empirycznie udowodniony liniowy rząd zbieżności ($O(h)$).
\end{itemize}

\end{document}
